\chapter{Ameaças à validade}
%não se esqueça de definir uma label única para utilizar no comando \ref
\label{chap:ameacas_validade}

Imagine que seu TCC é um castelo. Um castelo feito de fonte Arial 12, gráficos de pizza (vish), tabelas e aquele parágrafo lindo que você copiou... digo, se inspirou de um artigo gringo. Só que como toda construção, eventualmente ficam aparentes algumas rachaduras. Essas rachaduras são as ameaças à validade, que são basicamente tudo que pode fazer o seu trabalho parecer menos confiável. É aquela hora que chega alguém e diz: “Legal seu experimento... Mas você testou isso em três pessoas e um gato?”

Dessa forma, antecipar-se a essas ameaças é uma ótima forma de demonstrar que você já levou essas possibilidades em consideração, e não está alheio a esses riscos. Parte desse texto pode servir como insumo para a seção de trabalhos futuros.